\documentclass[11pt]{article}
\usepackage[margin=1in]{geometry}
\usepackage{booktabs}
\usepackage{longtable}
\title{Replacement memo: corrected heavy-tail Section 5.2}
\author{SVP simulation audit}
\date{15 September 2026}
\begin{document}
\maketitle

\section*{Decision}
Section 5.2 is replaced by the Student-$t(2)$ stress study described below.
The original \texttt{SVP\_NEW.pdf} is unchanged because no editable source was
available. The previous robust outputs are preserved in
\texttt{simulations/other\_simus/power\_robust\_legacy\_20260915/}.

\section*{Calibration}
RFPOP uses the biweight loss (called \texttt{Outlier} by \texttt{robseg}) with
$\ell$-threshold 3 and $\lambda=c_{\rm RF}\log n$. The smallest candidate with
corrected null F1 at least 0.99 is selected. Median-Mood and Wilcoxon SVP
thresholds are then selected to match the RFPOP null F1 target. Calibration and
validation use independent t(2) streams with 20,000 and 10,000 series,
respectively.; Wilson intervals and deterministic seeds are reported in the
CSV table below.

\begin{table}[ht]
\centering
\caption{Validation of selected robust null thresholds.}
\begin{tabular}{lllllll}
\toprule
Method & Candidate & F1 & Lower & Upper & RFPOP target & Seed
 \\
\midrule
RFPOP & 3.0000 & 0.9958 & 0.9943 & 0.9969 & 0.9961 & 920000.0000 \\
Wilcoxon & 1.7500 & 0.9951 & 0.9935 & 0.9963 & 0.9961 & 920000.0000 \\
MedianMood & 0.0010 & 0.9979 & 0.9968 & 0.9986 & 0.9961 & 920000.0000 \\
\bottomrule
\end{tabular}
\end{table}


The corrected Gaussian-PELT null F1 is the external benchmark. Gaussian PELT
under t(2) is retained as a stress baseline and is not used as a calibration
target because its Gaussian SSE fit produces almost universal false positives.

\section*{Before/after metrics}
The detailed paired table is indexed by pattern and jump in
\texttt{paired\_comparison.csv}; the compact table below averages changes over
the 40 jump sizes within each pattern. Positive changes improve F1, precision,
recall and correct-number probability; negative changes improve MSE and segment
count.

\begin{table}[ht]
\centering
\caption{Mean new-minus-legacy metric change by pattern and method.}
\begin{tabular}{llllllll}
\toprule
Pattern & Method & F1 & Precision & Recall & Correct & MSE & Segments
 \\
\midrule
none & MedianMood & 0.0041 & 0.0069 & -0.0015 & 0.0153 & -0.0006 & -0.0155 \\
none & PELT & -0.0258 & -0.0131 & -1.0000 & 0.0000 & 0.0000 & 0.0000 \\
none & RFPOP & 0.0476 & 0.0666 & -0.0050 & 0.1105 & -0.0068 & -0.2020 \\
none & Wilcoxon & 0.0031 & 0.0066 & -0.0040 & 0.0173 & -0.0014 & -0.0173 \\
rand1 & MedianMood & -0.0418 & -0.1757 & -0.0019 & 0.2297 & -0.0895 & 0.8985 \\
rand1 & PELT & -0.0202 & -0.0117 & -0.0421 & 0.0000 & 0.0000 & 0.0000 \\
rand1 & RFPOP & -0.0690 & -0.0968 & -0.0686 & -0.0050 & 0.0008 & -0.2902 \\
rand1 & Wilcoxon & -0.0537 & -0.0510 & -0.0565 & -0.0365 & -0.0100 & -0.1928 \\
up & MedianMood & -0.0397 & -0.0910 & -0.0215 & 0.1145 & -0.1017 & 0.4695 \\
up & PELT & -0.0202 & -0.0117 & -0.0417 & 0.0000 & 0.0000 & 0.0000 \\
up & RFPOP & -0.1183 & -0.1049 & -0.1052 & 0.0000 & 0.0734 & -0.2760 \\
up & Wilcoxon & -0.0620 & -0.0634 & -0.0609 & -0.0285 & -0.0261 & -0.0975 \\
updown & MedianMood & -0.0376 & -0.2366 & -0.0057 & 0.1350 & -0.0609 & 0.8145 \\
updown & PELT & -0.0202 & -0.0117 & -0.0413 & 0.0000 & 0.0000 & 0.0000 \\
updown & RFPOP & -0.0858 & -0.1698 & -0.0746 & -0.0400 & 0.0060 & -0.6450 \\
updown & Wilcoxon & -0.0633 & -0.0647 & -0.0622 & -0.0312 & 0.0011 & -0.1077 \\
\bottomrule
\end{tabular}
\end{table}


\section*{Replacement Section 5.2}

We assess robustness to heavy-tailed noise using independent Student-$t(2)$
errors generated by \texttt{ts\_generator(type = "student", df = 2, scale = 1)}.
The design has $n=1000$, four patterns (none, up, updown and rand1), jump sizes
from 0.1 to 4 in increments of 0.1, and 100 replications. The fixed Gaussian
PELT fit with penalty $2\log n$ is retained as an intentionally misspecified
stress baseline. The robust competitors are null-calibrated RFPOP with
biweight loss and $\ell$-threshold 3, Median-Mood SVP, and Wilcoxon SVP.

For every method, the mandatory terminal endpoint $n$ is removed before
one-to-one matching within the framework tolerance. We report precision,
recall, F1, correct-number probability, signal MSE, number of fitted segments,
and localization error. RFPOP is calibrated to at least 0.99 corrected null F1.
The Median-Mood and Wilcoxon thresholds are selected to match that target and
all SVP calls use the \texttt{right} pruning mode. Their length-dependent
scaling uses $K$, passed in the fitting code as \texttt{oracle\_segments}, the
known number of data-generating segments in this simulation. This oracle is a
design device and is not assumed available in practical data analysis.

\[
 \gamma_W(c,K)=c\sqrt{\frac{(n/K)^3}{12}},\qquad
 \gamma_M(\alpha,K)=\chi^2_{1,1-\alpha_s},\quad
 \alpha_s=1-(1-\alpha)^{1/(n/K-1)}.
\]

\end{document}
